\documentclass[11pt,letterpaper]{article}
\usepackage[margin=0.7in]{geometry}
\usepackage[T1]{fontenc}
\usepackage{lmodern}
\usepackage{amsmath,amssymb,booktabs,array,xcolor,listings,hyperref,fancyhdr}
\definecolor{accent}{RGB}{25,65,100}
\hypersetup{colorlinks=true,linkcolor=accent,urlcolor=accent}
\setlength{\parindent}{0pt}
\setlength{\parskip}{4pt}
\setlength{\headheight}{14pt}
\pagestyle{fancy}\fancyhf{}\fancyhead[L]{\small CSI2110 -- Study Notes}\fancyhead[R]{\small Fall 2026}\fancyfoot[C]{\thepage}
\lstset{basicstyle=\small\ttfamily,breaklines=true,columns=fullflexible,keepspaces=true,frame=single,rulecolor=\color{accent!25},backgroundcolor=\color{accent!3},showstringspaces=false,tabsize=2}
\newcommand{\sourcepages}[1]{\textit{Source: PDF pages #1.}}
\newcommand{\keyidea}[1]{\par\textbf{Key idea.} #1\par}
\newcommand{\cost}[1]{\ensuremath{O(#1)}}
\newcommand{\method}[1]{\texttt{#1}}

\title{\textbf{Lecture 8: (2,4) trees}\\\large CSI2110 -- Data Structures and Algorithms}\author{Lecture study notes}\date{}\begin{document}\maketitle
\textbf{Source document:} \path{../lec05-(2,4) trees.pdf}. These notes synthesize the source and add worked explanations. References use PDF page numbers. Pseudocode includes boundary cases.
\tableofcontents\newpage

\section{Multi-Way Search Tree}
A multi-way search tree is an ordered tree such that:
\begin{itemize}
    \item Nodes can store more than one element/key.
    \subitem Children are denoted $v_1$, $v_2$, ... , $v_{d-1}$ 
    \subitem Elements are denoted $k_1$, $k_2$, ... , $k_{d-1}$
    \item All elements are sorted within the node, and children (and their elements) are sorted respective to their parent node's elements 
    \item Leaves store no entries, left for placeholding
\end{itemize}

\section{Multi-Way Searching}
\keyidea{How do we search for elements in a multi way tree?}
Given an element $k_i$, we traverse the nodes and their elements.
\begin{itemize}
    \item 
\end{itemize}

\end{document}
